%------------------------------------------------------------------------------
% 3.4 حالات الاستخدام
%------------------------------------------------------------------------------
\section{حالات الاستخدام}
\label{sec:use-cases}

توضح حالات الاستخدام كيفية تفاعل مستخدمي نظام «بيان» مع وظائفه الرئيسية، كما تحدد الحدود بين الفاعلين البشريين والمكونات الآلية التي تنفذ المعالجة الداخلية. ويعتمد هذا القسم على ثماني حالات استخدام تغطي الاستعلام عن البيانات، والبحث في المستندات، وإعداد التقارير، وإدارة المستخدمين ومصادر البيانات، والتدقيق، والتهيئة الأولية للنظام.

\subsection{نموذج الفاعلين}
\label{subsec:actor-model}

يضم النظام ثمانية فاعلين رئيسيين ينقسمون إلى فئتين: فاعلون خارجيون بشريون وفاعلون داخليون آليون. وقد صُمم هذا الفصل بحيث يتعامل المستخدم البشري مع نقطة دخول موحدة، بينما تتولى الوكلاء المتخصصون تحليل الطلب وتنفيذه وفق صلاحيات الوصول المحددة.

\begin{figure}[htbp]
  \centering
  % Actor classification — v3 (2026-09): two-level actor taxonomy.
% Level 1 = actor KIND (pattern): external human actor (top band,
% 4 types) vs internal AI-agent actor (bottom band, 5 types —
% Reporter Agent added for full consistency with ch1/ch2/arch).
% Every Arabic label is wrapped in \textarabic{} inside bidi's LTR
% environment (same proven fix as the Gantt charts) so word order
% renders correctly. Card order runs right-to-left to match Arabic
% reading; numbered badges (1-9) match the actors table in 3.4.1.
\newsavebox{\ActorsBox}%
\begin{lrbox}{\ActorsBox}%
\begin{LTR}%
\begin{tikzpicture}[
  font=\small,
  stickman/.pic={
    \draw[line width=1.1pt, #1] (0,0.4) circle (0.2);
    \draw[line width=1.1pt, #1] (0,0.2) -- (0,-0.3);
    \draw[line width=1.1pt, #1] (-0.3,0.06) -- (0,-0.06) -- (0.3,0.06);
    \draw[line width=1.1pt, #1] (0,-0.3) -- (-0.22,-0.75);
    \draw[line width=1.1pt, #1] (0,-0.3) -- (0.22,-0.75);
  },
  robot/.pic={
    \draw[line width=1pt, #1] (0,0.44) -- (0,0.58);
    \fill[#1] (0,0.65) circle (0.05);
    \draw[line width=1pt, #1, rounded corners=2.5pt] (-0.24,0.14) rectangle (0.24,0.44);
    \fill[#1] (-0.11,0.29) circle (0.045);
    \fill[#1] (0.11,0.29) circle (0.045);
    \draw[line width=1pt, #1, rounded corners=3pt] (-0.3,-0.44) rectangle (0.3,0.08);
    \draw[line width=1pt, #1] (0,-0.18) circle (0.1);
    \fill[#1] (0,-0.18) circle (0.035);
  },
  actorcard/.style={
    rounded corners=6pt, draw=#1!75, fill=white, line width=0.9pt,
    minimum width=3.7cm, minimum height=3.05cm
  },
  actornum/.style={
    circle, fill=#1, text=white, font=\bfseries\scriptsize,
    inner sep=0pt, minimum size=0.42cm
  },
  actorname/.style={
    align=center, font=\bfseries\scriptsize, text=kaudark,
    inner sep=0pt
  },
  actoreng/.style={
    align=center, font=\tiny, text=kaudark!70, inner sep=0pt
  },
  actordesc/.style={
    align=center, font=\tiny, text=kaudark!85, inner sep=0pt,
    text width=3.25cm
  },
  actordescn/.style={
    align=center, font=\tiny, text=kaudark!85, inner sep=0pt,
    text width=2.5cm
  },
  bandpill/.style={
    rounded corners=4pt, fill=#1, text=white, inner sep=5pt,
    font=\bfseries\scriptsize
  }
]

  % ================= KIND 1: external human actor (4 types) =================
  \fill[kauprimary!6, rounded corners=10pt] (0.15,0.62) rectangle (15.85,-3.95);
  \node[bandpill=kauprimary] at (8.0,0.55)
    {\textarabic{النمط الأول: الفاعل البشري الخارجي — أربعة أنواع}};

  % Type 1 (rightmost): Enterprise Employee
  \draw[actorcard=kauprimary] (12.1,-0.55) rectangle (15.8,-3.6);
  \pic at (13.95,-1.35) {stickman={kaudark}};
  \node[actornum=kauprimary] at (15.45,-0.85) {\EN{1}};
  \node[actorname] at (13.95,-2.28) {\textarabic{موظف المؤسسة}};
  \node[actoreng]  at (13.95,-2.62) {\EN{Enterprise Employee}};
  \node[actordesc] at (13.95,-3.05) {\textarabic{إدخال الاستعلامات الأساسية والبحث في اللوائح والسياسات}};

  % Type 2: Enterprise Manager
  \draw[actorcard=kauprimary] (8.15,-0.55) rectangle (11.85,-3.6);
  \pic at (10.0,-1.35) {stickman={kaudark}};
  \node[actornum=kauprimary] at (11.5,-0.85) {\EN{2}};
  \node[actorname] at (10.0,-2.28) {\textarabic{مدير المؤسسة}};
  \node[actoreng]  at (10.0,-2.62) {\EN{Enterprise Manager}};
  \node[actordesc] at (10.0,-3.05) {\textarabic{الاطلاع على التقارير المتقدمة ومتابعة المؤشرات التشغيلية}};

  % Type 3: System Administrator
  \draw[actorcard=kauprimary] (4.2,-0.55) rectangle (7.9,-3.6);
  \pic at (6.05,-1.35) {stickman={kaudark}};
  \node[actornum=kauprimary] at (7.55,-0.85) {\EN{3}};
  \node[actorname] at (6.05,-2.28) {\textarabic{مسؤول النظام}};
  \node[actoreng]  at (6.05,-2.62) {\EN{System Administrator}};
  \node[actordesc] at (6.05,-3.05) {\textarabic{إدارة المستخدمين والأدوار وضبط الصلاحيات وسجلات التدقيق}};

  % Type 4: Onboarding Operator
  \draw[actorcard=kauprimary] (0.25,-0.55) rectangle (3.95,-3.6);
  \pic at (2.1,-1.35) {stickman={kaudark}};
  \node[actornum=kauprimary] at (3.6,-0.85) {\EN{4}};
  \node[actorname] at (2.1,-2.28) {\textarabic{مشغل التهيئة}};
  \node[actoreng]  at (2.1,-2.62) {\EN{Onboarding Operator}};
  \node[actordesc] at (2.1,-3.05) {\textarabic{ربط النظام بالمؤسسات الجديدة وتهيئة البيانات لأول مرة}};

  % ============ KIND 2: internal AI-agent actor (5 types) ============
  \fill[kausecondary!6, rounded corners=10pt] (0.15,-4.35) rectangle (15.85,-8.92);
  \node[bandpill=kausecondary] at (8.0,-4.42)
    {\textarabic{النمط الثاني: الفاعل الداخلي وكيل الذكاء الاصطناعي — خمسة أنواع}};

  % Type 5 (rightmost): Supervisor Agent — unified entry point
  \draw[actorcard=kausecondary] (12.88,-5.52) rectangle (15.8,-8.57);
  \pic at (14.34,-6.32) {robot={kausecondary!90!black}};
  \node[actornum=kausecondary] at (15.45,-5.82) {\EN{5}};
  \node[actorname] at (14.34,-7.25) {\textarabic{الوكيل المشرف}};
  \node[actoreng]  at (14.34,-7.59) {\EN{Supervisor Agent}};
  \node[actordescn] at (14.34,-8.02) {\textarabic{نقطة الدخول الموحدة وتحليل النية وتوجيه الطلبات}};

  % Type 6: SQL Agent
  \draw[actorcard=kausecondary] (9.71,-5.52) rectangle (12.63,-8.57);
  \pic at (11.17,-6.32) {robot={kausecondary!90!black}};
  \node[actornum=kausecondary] at (12.28,-5.82) {\EN{6}};
  \node[actorname] at (11.17,-7.25) {\textarabic{وكيل قواعد البيانات}};
  \node[actoreng]  at (11.17,-7.59) {\EN{SQL Agent}};
  \node[actordescn] at (11.17,-8.02) {\textarabic{تحويل الأسئلة العربية إلى استعلامات قواعد البيانات}};

  % Type 7: RAG Agent
  \draw[actorcard=kausecondary] (6.54,-5.52) rectangle (9.46,-8.57);
  \pic at (8.0,-6.32) {robot={kausecondary!90!black}};
  \node[actornum=kausecondary] at (9.11,-5.82) {\EN{7}};
  \node[actorname] at (8.0,-7.25) {\textarabic{وكيل الاسترجاع}};
  \node[actoreng]  at (8.0,-7.59) {\EN{RAG Agent}};
  \node[actordescn] at (8.0,-8.02) {\textarabic{استرجاع المقاطع النصية من المستندات واللوائح}};

  % Type 8: Analyst Agent
  \draw[actorcard=kausecondary] (3.37,-5.52) rectangle (6.29,-8.57);
  \pic at (4.83,-6.32) {robot={kausecondary!90!black}};
  \node[actornum=kausecondary] at (5.94,-5.82) {\EN{8}};
  \node[actorname] at (4.83,-7.25) {\textarabic{وكيل التحليل}};
  \node[actoreng]  at (4.83,-7.59) {\EN{Analyst Agent}};
  \node[actordescn] at (4.83,-8.02) {\textarabic{التحليل الإحصائي واكتشاف الانحرافات والقيم الشاذة}};

  % Type 9: Reporter Agent
  \draw[actorcard=kausecondary] (0.2,-5.52) rectangle (3.12,-8.57);
  \pic at (1.66,-6.32) {robot={kausecondary!90!black}};
  \node[actornum=kausecondary] at (2.77,-5.82) {\EN{9}};
  \node[actorname] at (1.66,-7.25) {\textarabic{وكيل التقارير}};
  \node[actoreng]  at (1.66,-7.59) {\EN{Reporter Agent}};
  \node[actordescn] at (1.66,-8.02) {\textarabic{تنسيق المخرجات وتوليد التقارير المنسقة للتصدير}};

\end{tikzpicture}%
\end{LTR}%
\end{lrbox}%
\resizebox{\textwidth}{!}{\usebox{\ActorsBox}}%

  \caption{تصنيف فاعلي نظام «بيان» إلى فاعلين بشريين ووكلاء آليين.}
  \label{fig:use-case-actors}
\end{figure}

\begin{table}[htbp]
\centering
\caption{الفاعلون الرئيسيون في نظام «بيان».}
\label{tab:use-case-actors}
\small
\renewcommand{\arraystretch}{1.5}
\begin{tabularx}{\textwidth}{>{\raggedright\arraybackslash}p{3.2cm} >{\raggedright\arraybackslash}p{3.1cm} X}
\toprule
\textbf{الفئة} & \textbf{الفاعل} & \textbf{المسؤوليات والتفاعلات الرئيسية} \\
\midrule
\multirow{4}{*}{فاعل خارجي بشري}
& موظف المؤسسة \EN{(Enterprise Employee)} & إدخال الاستعلامات الأساسية والبحث في اللوائح والسياسات الخاصة بعمله. \\
\cmidrule{2-3}
& مدير المؤسسة \EN{(Enterprise Manager)} & الاطلاع على التعديلات والتقارير المتقدمة، ومتابعة المؤشرات التشغيلية والمالية. \\
\cmidrule{2-3}
& مسؤول النظام \EN{(System Administrator)} & إدارة المستخدمين والأدوار، ضبط الصلاحيات، متابعة سجلات التدقيق وإدارة مصادر البيانات. \\
\cmidrule{2-3}
& مشغّل التهيئة \EN{(Onboarding Operator)} & المسؤول عن ربط النظام بالمؤسسات الجديدة وإعداد قواعد البيانات والمستندات لأول مرة. \\
\midrule
\multirow{4}{*}{فاعل داخلي آلي}
& الوكيل المشرف \EN{(Supervisor Agent)} & يمثل نقطة الدخول الموحدة التي تتلقى طلبات المستخدم، تحلل النية \EN{(Intent Classification)}، وتوجه الطلب للوكيل المختص. \\
\cmidrule{2-3}
& وكيل \EN{SQL} & تحويل الأسئلة العربية إلى استعلامات \EN{SQL} ومعالجتها على قواعد البيانات. \\
\cmidrule{2-3}
& وكيل \EN{RAG} & استرجاع المقاطع والدلائل النصية من المستندات واللوائح غير المهيكلة. \\
\cmidrule{2-3}
& وكيل المحلل \EN{(Analyst Agent)} & معالجة البيانات الإحصائية، اكتشاف الانحرافات والقيم الشاذة، وتوليد الرسوم البيانية والتقارير. \\
\bottomrule
\end{tabularx}
\end{table}

\begin{figure}[htbp]
  \centering
  % Use Case Diagram — redesigned v2 (2026-09).
% Fixes Arabic word reversal by wrapping the picture in bidi's LTR
% environment and every Arabic label in \textarabic{} (same proven
% pattern as the chapter-1 Gantt charts).
% Layout: human actors on the RIGHT (matches section 3.4.2 prose),
% AI agents on the LEFT, 8 use cases in a single central column
% split into two tinted zones (user services / administration).
% The internal "Route Query" hub connects to the supervisor agent
% (execution line), which delegates to the three specialist agents
% through a vertical delegation bus. include/extend relations are
% routed through the empty side corridors so no line crosses any
% shape, and every label sits clear of lines and edges.
\newsavebox{\UCBox}%
\begin{lrbox}{\UCBox}%
\begin{LTR}%
\begin{tikzpicture}[
  font=\small,
  stickman/.pic={
    \draw[line width=1.1pt, #1] (0,0.4) circle (0.2);
    \draw[line width=1.1pt, #1] (0,0.2) -- (0,-0.3);
    \draw[line width=1.1pt, #1] (-0.3,0.06) -- (0,-0.06) -- (0.3,0.06);
    \draw[line width=1.1pt, #1] (0,-0.3) -- (-0.22,-0.75);
    \draw[line width=1.1pt, #1] (0,-0.3) -- (0.22,-0.75);
  },
  robot/.pic={
    \draw[line width=1pt, #1] (0,0.44) -- (0,0.58);
    \fill[#1] (0,0.65) circle (0.05);
    \draw[line width=1pt, #1, rounded corners=2.5pt] (-0.24,0.14) rectangle (0.24,0.44);
    \fill[#1] (-0.11,0.29) circle (0.045);
    \fill[#1] (0.11,0.29) circle (0.045);
    \draw[line width=1pt, #1, rounded corners=3pt] (-0.3,-0.44) rectangle (0.3,0.08);
    \draw[line width=1pt, #1] (0,-0.18) circle (0.1);
    \fill[#1] (0,-0.18) circle (0.035);
  },
  ucellipse/.style={
    draw=kauprimary!70, fill=kauprimary!14, line width=0.8pt,
    ellipse, align=center, inner xsep=6pt, inner ysep=4pt,
    text=kaudark, font=\bfseries\scriptsize,
    minimum width=4.0cm, minimum height=1.15cm
  },
  ucadmin/.style={
    ucellipse, draw=kausecondary!70, fill=kausecondary!13
  },
  ucroute/.style={
    ellipse, draw=kauaccent, fill=kauaccent!8, dashed,
    line width=0.9pt, align=center, inner xsep=5pt, inner ysep=3pt,
    text=kauaccent!90!black, font=\bfseries\scriptsize
  },
  humanchip/.style={
    rounded corners=3pt, draw=kaudark!50, fill=white,
    align=center, inner sep=4pt, text=kaudark,
    font=\scriptsize\bfseries, minimum width=2.35cm
  },
  agentchip/.style={
    rounded corners=3pt, draw=kausecondary!60, fill=white,
    align=center, inner sep=3.5pt, text=kausecondary!90!black,
    font=\scriptsize\bfseries, minimum width=2.0cm
  },
  assoc/.style={draw=kaudark!75, line width=0.8pt},
  dep/.style={draw=kauaccent!90, line width=0.8pt, dashed, -stealth},
  exec/.style={draw=kausecondary!90, line width=0.9pt},
  deleg/.style={draw=kausecondary!75, line width=0.7pt}
]

  % ---- Zone tints (behind everything) ----
  \fill[kauprimary!6, rounded corners=9pt] (3.35,-1.42) rectangle (10.55,-7.42);
  \fill[kausecondary!6, rounded corners=9pt] (3.35,-7.62) rectangle (10.55,-13.54);
  \draw[kauprimary!45, dotted, line width=0.6pt] (3.35,-7.52) -- (10.55,-7.52);

  % ---- System boundary + title tab ----
  \draw[kauprimary!80, line width=1.1pt, rounded corners=10pt]
    (3.1,1.15) rectangle (10.8,-13.75);
  \node[fill=kauprimary, text=white, rounded corners=4pt, inner sep=6pt,
        font=\bfseries\small] at (6.95,1.15)
    {\textarabic{نظام بيان الذكي للاستعلام والتحليل}};

  % ---- Internal hub: Route Query ----
  \node[ucroute, minimum width=3.3cm, minimum height=1.02cm] (route) at (6.95,-0.68)
    {\textarabic{توجيه الاستعلام}\\[1pt]{\tiny\EN{Route Query (internal)}}};

  % ---- Use cases: user zone ----
  \node[ucellipse] (uc1) at (6.95,-2.25) {\EN{UC-01}\\[1pt]\textarabic{طرح استعلام عن البيانات}};
  \node[ucellipse] (uc2) at (6.95,-3.73) {\EN{UC-02}\\[1pt]\textarabic{البحث في المستندات}};
  \node[ucellipse] (uc3) at (6.95,-5.21) {\EN{UC-03}\\[1pt]\textarabic{توليد تقرير مخصص}};
  \node[ucellipse] (uc4) at (6.95,-6.69) {\EN{UC-04}\\[1pt]\textarabic{عرض الإشعارات}};

  % ---- Use cases: administration zone ----
  \node[ucadmin] (uc5) at (6.95,-8.37) {\EN{UC-05}\\[1pt]\textarabic{إدارة المستخدمين والأدوار}};
  \node[ucadmin] (uc6) at (6.95,-9.85) {\EN{UC-06}\\[1pt]\textarabic{إدارة مصادر البيانات}};
  \node[ucadmin] (uc7) at (6.95,-11.33) {\EN{UC-07}\\[1pt]\textarabic{مراجعة سجل التدقيق}};
  \node[ucadmin] (uc8) at (6.95,-12.81) {\EN{UC-08}\\[1pt]\textarabic{تشغيل التهيئة الذكية}};

  % ---- HUMAN ACTORS (right side) ----
  \pic at (14.35,-3.45) {stickman={kaudark}};
  \node[humanchip, below=3pt] at (14.35,-4.42)
    {\textarabic{موظف المؤسسة}\\[1pt]{\tiny\EN{Enterprise Employee}}};

  \pic at (14.35,-6.10) {stickman={kaudark}};
  \node[humanchip, below=3pt] at (14.35,-7.07)
    {\textarabic{مدير المؤسسة}\\[1pt]{\tiny\EN{Enterprise Manager}}};

  \pic at (14.35,-9.30) {stickman={kaudark}};
  \node[humanchip, below=3pt] at (14.35,-10.27)
    {\textarabic{مسؤول النظام}\\[1pt]{\tiny\EN{System Administrator}}};

  \pic at (14.35,-12.40) {stickman={kaudark}};
  \node[humanchip, below=3pt] at (14.35,-13.37)
    {\textarabic{مشغل التهيئة}\\[1pt]{\tiny\EN{Onboarding Operator}}};

  % ---- AI AGENT ACTORS (left side) ----
  \pic at (1.45,-1.75) {robot={kausecondary!90!black}};
  \node[agentchip, below=3pt] at (1.45,-2.72)
    {\textarabic{الوكيل المشرف}\\[1pt]{\tiny\EN{Supervisor Agent}}};

  \pic at (1.45,-4.30) {robot={kausecondary!90!black}};
  \node[agentchip, below=3pt] at (1.45,-5.27)
    {\textarabic{وكيل قواعد البيانات}\\[1pt]{\tiny\EN{SQL Agent}}};

  \pic at (1.45,-6.90) {robot={kausecondary!90!black}};
  \node[agentchip, below=3pt] at (1.45,-7.87)
    {\textarabic{وكيل الاسترجاع}\\[1pt]{\tiny\EN{RAG Agent}}};

  \pic at (1.45,-9.40) {robot={kausecondary!90!black}};
  \node[agentchip, below=3pt] at (1.45,-10.37)
    {\textarabic{وكيل التحليل}\\[1pt]{\tiny\EN{Analyst Agent}}};

  % ---- Associations: employee (UC-01..04) ----
  \draw[assoc] (14.0,-3.45) -- (uc1.east);
  \draw[assoc] (14.0,-3.45) -- (uc2.east);
  \draw[assoc] (14.0,-3.45) -- (uc3.east);
  \draw[assoc] (14.0,-3.45) -- (uc4.east);

  % ---- Associations: manager (UC-01..04, offset anchors) ----
  \draw[assoc] (14.0,-6.10) -- (uc1.south east);
  \draw[assoc] (14.0,-6.10) -- (uc2.south east);
  \draw[assoc] (14.0,-6.10) -- (uc3.south east);
  \draw[assoc] (14.0,-6.10) -- (uc4.north east);

  % ---- Associations: system administrator (UC-04..07) ----
  \draw[assoc] (14.0,-9.30) -- (uc4.south east);
  \draw[assoc] (14.0,-9.30) -- (uc5.north east);
  \draw[assoc] (14.0,-9.30) -- (uc6.north east);
  \draw[assoc] (14.0,-9.30) -- (uc7.north east);

  % ---- Associations: onboarding operator (UC-08) ----
  \draw[assoc] (14.0,-12.40) -- (uc8.north east);

  % ---- Include / extend (routed through empty side corridors) ----
  \draw[dep] (uc1.north east) -- (route.east)
    node[midway, above right=-1pt, font=\scriptsize\itshape, text=kaudark!85,
         fill=white, inner sep=1.5pt] {\EN{<<include>>}};
  \draw[dep] (uc2.north west)
    .. controls (4.30,-2.80) and (4.30,-1.25) .. (route.west)
    node[pos=0.30, font=\scriptsize\itshape, text=kaudark!85,
         fill=white, inner sep=1.5pt] {\EN{<<include>>}};
  \draw[dep] (uc3.west)
    .. controls (4.75,-4.90) and (4.75,-2.60) .. (uc1.west)
    node[pos=0.50, xshift=-0.35cm, font=\scriptsize\itshape, text=kaudark!85,
         fill=white, inner sep=1.5pt] {\EN{<<extend>>}};

  % ---- Execution: hub -> supervisor ----
  \draw[exec] (route.south west) -- (2.90,-1.75)
    node[pos=0.62, anchor=south, font=\scriptsize\bfseries, text=kausecondary!90!black,
         inner sep=1.5pt] {\textarabic{تنفيذ}};

  % ---- Delegation bus: supervisor -> specialist agents ----
  \draw[deleg] (2.90,-1.75) -- (2.90,-9.40);
  \draw[deleg] (2.90,-1.75) -- (1.80,-1.75);
  \draw[deleg, -{stealth[scale=0.75]}] (2.90,-4.30) -- (1.85,-4.30);
  \draw[deleg, -{stealth[scale=0.75]}] (2.90,-6.90) -- (1.85,-6.90);
  \draw[deleg, -{stealth[scale=0.75]}] (2.90,-9.40) -- (1.85,-9.40);

  % ---- Legend ----
  \node[draw=dividerdim!80, fill=white, rounded corners=3pt, inner sep=6pt,
        anchor=north west, font=\tiny, text=kaudark] at (0.10,-11.30)
  {%
    \begin{tabular}{@{}c@{\hspace{5pt}}r@{}}
      \tikz{\pic[scale=0.62, yshift=0.02] {stickman={kaudark}};} & \textarabic{فاعل بشري خارجي} \\[3pt]
      \tikz{\pic[scale=0.62] {robot={kausecondary!90!black}};} & \textarabic{وكيل ذكاء اصطناعي} \\[3pt]
      \tikz{\fill[kauprimary!14, draw=kauprimary!70] (0.2,0.1) ellipse [x radius=0.17, y radius=0.095];} & \textarabic{حالة استخدام} \\[3pt]
      \tikz{\draw[assoc] (0,0.1) -- (0.42,0.1);} & \textarabic{ارتباط فاعل بحالة} \\[3pt]
      \tikz{\draw[dep] (0,0.1) -- (0.42,0.1);} & \textarabic{علاقة تضمين أو توسيع} \\[3pt]
      \tikz{\draw[deleg] (0.02,0.18) -- (0.02,0.02) -- (0.42,0.02);} & \textarabic{تفويض بين الوكلاء} \\
    \end{tabular}%
  };

\end{tikzpicture}%
\end{LTR}%
\end{lrbox}%
\resizebox{\textwidth}{!}{\usebox{\UCBox}}%

  \caption{مخطط حالات استخدام نظام «بيان».}
  \label{fig:use-case-diagram}
\end{figure}

\subsection{نطاق حالات الاستخدام}
\label{subsec:use-case-scope}

تغطي الحالات الثماني الوحدات الوظيفية الأساسية للنظام، كما يوضح الجدول \ref{tab:use-case-summary}. وتُعد الحالات ذات الأولوية العالية جزءاً من الإصدار الأساسي، في حين تمثل الإشعارات وظيفة تحسين يمكن تنفيذها بعد استقرار الوظائف الجوهرية.

\begin{table}[htbp]
\centering
\caption{الملخص المرجعي لحالات استخدام نظام «بيان».}
\label{tab:use-case-summary}
\small
\renewcommand{\arraystretch}{1.45}
\begin{tabularx}{\textwidth}{>{\centering\arraybackslash}p{1.3cm} >{\raggedright\arraybackslash}p{4.3cm} >{\raggedright\arraybackslash}p{4.3cm} >{\centering\arraybackslash}p{1.8cm}}
\toprule
\textbf{المعرّف} & \textbf{حالة الاستخدام} & \textbf{الفاعل الرئيسي} & \textbf{الأولوية} \\
\midrule
UC-01 & طرح استعلام عن البيانات & موظف المؤسسة / مدير المؤسسة & عالية \\
\midrule
UC-02 & البحث في المستندات واللوائح & موظف المؤسسة / مدير المؤسسة & عالية \\
\midrule
UC-03 & توليد تقرير مخصص & مدير المؤسسة / موظف المؤسسة & عالية \\
\midrule
UC-04 & عرض الإشعارات والتنبيهات & موظف المؤسسة / مدير المؤسسة / مسؤول النظام & منخفضة \\
\midrule
UC-05 & إدارة المستخدمين والأدوار & مسؤول النظام & عالية \\
\midrule
UC-06 & إدارة مصادر البيانات & مسؤول النظام & عالية \\
\midrule
UC-07 & مراجعة سجل التدقيق & مسؤول النظام & عالية \\
\midrule
UC-08 & تشغيل التهيئة الذكية & مشغل التهيئة & عالية \\
\bottomrule
\end{tabularx}
\end{table}

\begin{figure}[htbp]
  \centering
  % Sequence Diagram – faithful to original with improved formatting.
% 9 Lifelines: User, Frontend, Supervisor, SQL Agent, RAG Agent, Analyst Agent, PostgreSQL, Qdrant, Audit Log.
\resizebox{\textwidth}{!}{%
\begin{tikzpicture}[
  font=\scriptsize,
  header/.style={
    draw=kauprimary!60, fill=kauprimary!12, line width=0.9pt,
    rounded corners=3pt, align=center, inner sep=4pt,
    text=kaudark, font=\bfseries\scriptsize, minimum width=1.7cm, minimum height=0.7cm
  },
  lifeline/.style={draw=dividerdim, line width=0.6pt, dashed},
  msg/.style={-stealth, line width=0.7pt, draw=kaudark},
  reply/.style={-stealth, line width=0.7pt, draw=kaudark, dashed},
  mlabel/.style={font=\scriptsize, text=kaudark, fill=white, inner sep=1.5pt},
  altbox/.style={draw=kauprimary!50, line width=0.8pt, rounded corners=0pt, fill=white},
  altlabel/.style={fill=kauprimary!20, font=\bfseries\scriptsize, text=kaudark, inner sep=3pt},
  extbar/.style={fill=kauprimary!25, draw=kauprimary!50, line width=0.8pt},
  selfloop/.style={-stealth, line width=0.7pt, draw=kaudark},
  auditbar/.style={fill=kauprimary!30, draw=kauprimary!60, line width=0.8pt}
]

  % ---- X positions for 9 lifelines (right to left for RTL feel) ----
  \def\xUser{0}
  \def\xFront{2.6}
  \def\xSup{5.2}
  \def\xSql{7.8}
  \def\xRag{10.4}
  \def\xAna{13.0}
  \def\xPg{15.6}
  \def\xQd{18.2}
  \def\xAudit{20.8}

  % ---- Header boxes ----
  \node[header] (hUser) at (\xUser, 0) {المستخدم المؤسسي};
  \node[header] (hFront) at (\xFront, 0) {الواجهة الأمامية};
  \node[header] (hSup) at (\xSup, 0) {الوكيل المشرف};
  \node[header] (hSql) at (\xSql, 0) {وكيل \EN{SQL}};
  \node[header] (hRag) at (\xRag, 0) {وكيل \EN{RAG}};
  \node[header] (hAna) at (\xAna, 0) {وكيل المحلل};
  \node[header] (hPg) at (\xPg, 0) {\EN{PostgreSQL}};
  \node[header] (hQd) at (\xQd, 0) {\EN{Qdrant}};
  \node[header] (hAudit) at (\xAudit, 0) {سجل التدقيق};

  % ---- Lifelines ----
  \def\yEnd{-32.5}
  \foreach \x in {\xUser, \xFront, \xSup, \xSql, \xRag, \xAna, \xPg, \xQd, \xAudit} {
    \draw[lifeline] (\x, -0.5) -- (\x, \yEnd);
  }

  % ---- Macros ----
  \newcommand{\smsg}[4]{\draw[msg] (#1, #3) -- (#2, #3) node[midway, mlabel, above] {#4};}
  \newcommand{\rmsg}[4]{\draw[reply] (#1, #3) -- (#2, #3) node[midway, mlabel, above] {#4};}

  % ===========================================================================
  %  PHASE 1: UC-01 Query Data (SQL path)
  % ===========================================================================

  % User asks
  \smsg{\xUser}{\xFront}{-1.2}{يكتب سؤالاً (مثال: «مبيعات الربع الأخير»)}
  \smsg{\xFront}{\xSup}{-2.0}{يرسل السؤال مع سياق الجلسة}

  % Include: Intent Classification
  \draw[extbar] (\xSup - 0.85, -2.6) rectangle (\xSup + 0.85, -3.4);
  \node[font=\bfseries\tiny, text=kaudark] at (\xSup, -2.85)
    {توجيه الاستعلام \EN{<<include>>}};
  \node[font=\tiny, text=kaudark] at (\xSup, -3.15)
    {تحليل النية \EN{(Intent Classification)}};

  % ALT box
  \draw[altbox] (-0.9, -3.8) rectangle (21.7, -10.4);
  \node[altlabel, anchor=north west] at (-0.9, -3.8) {\EN{alt}};

  % SQL Path
  \smsg{\xSup}{\xSql}{-4.4}{يوجه الطلب إلى وكيل \EN{SQL}}

  % Self-loop on SQL agent
  \draw[selfloop] (\xSql + 0.15, -5.0) -- (\xSql + 0.9, -5.0) -- (\xSql + 0.9, -5.5) -- (\xSql + 0.15, -5.5);
  \node[mlabel, right] at (\xSql + 0.95, -5.25) {تحويل \EN{NL2SQL} + تطبيق \EN{Smart Views}};

  \smsg{\xSql}{\xPg}{-6.1}{تنفيذ استعلام \EN{SQL} (مع \EN{RBAC})}
  \rmsg{\xPg}{\xSql}{-6.8}{إرجاع النتيجة الخام}
  \rmsg{\xSql}{\xSup}{-7.5}{إرجاع الإجابة المنسقة (جدول/نص)}
  \rmsg{\xSup}{\xFront}{-8.2}{عرض النتيجة}
  \rmsg{\xFront}{\xUser}{-8.9}{عرض النتيجة الأولية}

  % ---- Divider: <<extend>> UC-03 ----
  \fill[extbar] (-0.9, -10.6) rectangle (21.7, -11.2);
  \node[font=\bfseries\scriptsize, text=kaudark] at (10.4, -10.9)
    {\EN{<<extend>>  UC-03} (توليد تقرير مخصص)};

  % UC-03 flow
  \smsg{\xUser}{\xFront}{-11.8}{يطلب تحويل النتيجة إلى تقرير تحليلي}
  \smsg{\xFront}{\xSup}{-12.5}{طلب إنشاء تقرير}
  \smsg{\xSup}{\xAna}{-13.2}{يوجه الطلب إلى وكيل المحلل}

  % Self-loop on Analyst
  \draw[selfloop] (\xAna + 0.15, -13.8) -- (\xAna + 0.9, -13.8) -- (\xAna + 0.9, -14.3) -- (\xAna + 0.15, -14.3);
  \node[mlabel, right] at (\xAna + 0.95, -14.05) {استخراج المؤشرات والاتجاهات};

  \draw[selfloop] (\xAna + 0.15, -14.8) -- (\xAna + 0.9, -14.8) -- (\xAna + 0.9, -15.3) -- (\xAna + 0.15, -15.3);
  \node[mlabel, right] at (\xAna + 0.95, -15.05) {تنسيق تقرير (نص/جدول/رسم بياني)};

  \rmsg{\xAna}{\xSup}{-15.9}{إرجاع التقرير الجاهز}
  \rmsg{\xSup}{\xFront}{-16.6}{عرض التقرير}
  \rmsg{\xFront}{\xUser}{-17.3}{عرض التقرير مع خيارات التصدير}

  % ===========================================================================
  %  PHASE 2: UC-02 Search Documents (RAG path)
  % ===========================================================================

  % divider
  \draw[dividerdim, line width=0.4pt, dashed] (-0.9, -18.0) -- (21.7, -18.0);

  \smsg{\xSup}{\xRag}{-18.8}{يوجه الطلب إلى وكيل \EN{RAG}}

  % Self-loop on RAG
  \draw[selfloop] (\xRag + 0.15, -19.4) -- (\xRag + 0.9, -19.4) -- (\xRag + 0.9, -19.9) -- (\xRag + 0.15, -19.9);
  \node[mlabel, right] at (\xRag + 0.95, -19.65) {توليد متجه البحث \EN{(Query Embedding)}};

  \smsg{\xRag}{\xQd}{-20.5}{استرجاع المقاطع الأقرب \EN{(Top-K)}};
  \rmsg{\xQd}{\xRag}{-21.2}{إرجاع المقاطع + المصادر}

  % Self-loop on RAG for answer generation
  \draw[selfloop] (\xRag + 0.15, -21.8) -- (\xRag + 0.9, -21.8) -- (\xRag + 0.9, -22.3) -- (\xRag + 0.15, -22.3);
  \node[mlabel, right] at (\xRag + 0.95, -22.05) {توليد الإجابة مع الاستشهاد};

  \rmsg{\xRag}{\xSup}{-22.9}{إرجاع الإجابة مع المصادر}
  \rmsg{\xSup}{\xFront}{-23.6}{عرض الإجابة}
  \rmsg{\xFront}{\xUser}{-24.3}{عرض الإجابة (نص + رابط المصدر)}

  % ===========================================================================
  %  AUDIT LOGGING
  % ===========================================================================

  \fill[auditbar] (-0.9, -25.2) rectangle (21.7, -25.9);
  \node[font=\bfseries\scriptsize, text=kaudark] at (10.4, -25.4)
    {\EN{(Agent Trace)} يتم تسجيل كامل التتبع};
  \node[font=\scriptsize, text=kaudark] at (10.4, -25.7)
    {في سجل التدقيق لأغراض الرقابة};

  \smsg{\xSup}{\xAudit}{-26.5}{تسجيل (المستخدم، السؤال، الوكيل، الإجابة، الزمن)}

  % ---- Footer: repeated headers ----
  \foreach \x/\lbl in {
    \xUser/{المستخدم المؤسسي},
    \xFront/{الواجهة الأمامية},
    \xSup/{الوكيل المشرف},
    \xSql/{وكيل \EN{SQL}},
    \xRag/{وكيل \EN{RAG}},
    \xAna/{وكيل المحلل},
    \xPg/{\EN{PostgreSQL}},
    \xQd/{\EN{Qdrant}},
    \xAudit/{سجل التدقيق}%
  } {
    \node[header] at (\x, \yEnd - 0.5) {\lbl};
  }

\end{tikzpicture}%
}

  \caption{مخطط التسلسل لتدفقات الاستعلام والبحث وتوليد التقرير.}
  \label{fig:use-case-sequence}
\end{figure}

\subsection{الوصف التفصيلي لحالات الاستخدام}
\label{subsec:detailed-use-cases}

يعرض الوصف الآتي الهدف والفاعلين والشروط المسبقة والنتائج ومسار التنفيذ والتدفقات البديلة لكل حالة استخدام. وتساعد هذه الصياغة على تحويل المتطلبات إلى سيناريوهات قابلة للتصميم والاختبار.

\subsubsection{UC-01: طرح استعلام عن البيانات}
\label{subsubsec:uc01}

\begin{table}[htbp]
\centering
\caption{الوصف التفصيلي لحالة الاستخدام UC-01: طرح استعلام عن البيانات.}
\label{tab:uc01}
\small
\renewcommand{\arraystretch}{1.45}
\begin{tabularx}{\textwidth}{>{\raggedright\arraybackslash}p{3.1cm} X}
\toprule
\textbf{الحقل} & \textbf{الوصف} \\
\midrule
المعرّف والاسم & \EN{UC-01 — Ask Data Query} \\
\midrule
الهدف & تمكين موظف أو مدير المؤسسة من طرح استعلام بيانات باللغة العربية الطبيعية (فصحى أو عامية) والحصول على إجابة مستخلصة من قاعدة بيانات المؤسسة مع تطبيق قيود الأمان. \\
\midrule
الفاعل الرئيسي & موظف المؤسسة / مدير المؤسسة. \\
\midrule
الفاعلون المساعدون & الوكيل المشرف، ووكيل \EN{SQL}. \\
\midrule
المحفز & كتابة المستخدم سؤالاً في نافذة المحادثة والضغط على زر الإرسال. \\
\midrule
الشروط المسبقة & أن يكون المستخدم مسجَّل الدخول إلى النظام، وأن تكون بياناته مصنّفة ضمن مستوى أمان يتيح له الوصول للجدول المستهدف. \\
\midrule
الشروط اللاحقة & عرض الإجابة على شكل نص أو جدول أو رسم بياني، مع تسجيل الاستعلام والإجابة وسجل التتبع في سجل التدقيق. \\
\midrule
المسار الأساسي & 1) يطرح المستخدم سؤالاً بالعربية، مثل: «كم عدد الطلاب المسجلين في كلية الهندسة؟».\\ 2) يوجّه الوكيل المشرف الاستعلام إلى وكيل \EN{SQL}.\\ 3) يحوّل وكيل \EN{SQL} السؤال إلى استعلام \EN{SQL} عبر طبقة \EN{Smart Views}.\\ 4) يُنفَّذ الاستعلام على قاعدة البيانات ضمن حدود صلاحية المستخدم \EN{(RBAC)}.\\ 5) تُحوَّل النتيجة إلى إجابة نصية أو رسم بياني ويُعرض للمستخدم. \\
\midrule
التدفقات البديلة & \textbf{أخطاء نحوية:} في حال تعذّر توليد استعلام \EN{SQL} صحيح، يعيد النظام المحاولة بتقنية \EN{Few-Shot}، وإذا فشل يعرض رسالة توضح غموض السؤال ويطلب من المستخدم إعادة صياغته.\\ \textbf{انتهاك الصلاحية:} إذا حاول الوصول لجدول محظور، يرفض النظام ويُسجّل المحاولة ويعرض رسالة تعذّر الوصول. \\
\bottomrule
\end{tabularx}
\end{table}

\subsubsection{UC-02: البحث في المستندات واللوائح}
\label{subsubsec:uc02}

\begin{table}[htbp]
\centering
\caption{الوصف التفصيلي لحالة الاستخدام UC-02: البحث في المستندات.}
\label{tab:uc02}
\small
\renewcommand{\arraystretch}{1.45}
\begin{tabularx}{\textwidth}{>{\raggedright\arraybackslash}p{3.1cm} X}
\toprule
\textbf{الحقل} & \textbf{الوصف} \\
\midrule
المعرّف والاسم & \EN{UC-02 — Search Documents} \\
\midrule
الهدف & تمكين موظف أو مدير المؤسسة من طرح سؤال مرتبط بمستند أو لائحة مؤسسية غير مهيكلة، والحصول على إجابة موثّقة بالمصدر (اسم المستند ورقم الصفحة والمقطع المقتبس). \\
\midrule
الفاعل الرئيسي & موظف المؤسسة / مدير المؤسسة. \\
\midrule
الفاعلون المساعدون & الوكيل المشرف، ووكيل \EN{RAG}. \\
\midrule
المحفز & طرح المستخدم سؤالاً متعلقاً بمحتوى وثائقي وليس ببيانات جدولية، مثل: «ما هي سياسة الإجازات؟». \\
\midrule
الشروط المسبقة & أن تكون المستندات المؤسسية ذات الصلة قد فُهرست مسبقاً ضمن قاعدة البيانات المتجهة \EN{(Qdrant)}. \\
\midrule
الشروط اللاحقة & عرض إجابة نصية مرفقة باسم المستند المصدر ورقم الصفحة أو المقطع المستشهد به، وتحديث سجل التدقيق بمسار الاسترجاع. \\
\midrule
المسار الأساسي & 1) يطرح المستخدم سؤالاً نصياً.\\ 2) يوجّه الوكيل المشرف الاستعلام إلى وكيل \EN{RAG}.\\ 3) يسترجع وكيل \EN{RAG} أقرب المقاطع دلالياً من قاعدة المتجهات \EN{(Qdrant)}.\\ 4) يولّد النموذج اللغوي إجابة معزَّزة بالمقاطع المسترجعة مع الاستشهاد بمصدرها.\\ 5) يُعرض الجواب للمستخدم مع رابط أو إشارة للمستند الأصلي. \\
\midrule
التدفق البديل & \textbf{نقص المعلومات:} في حال عدم العثور على مقاطع ذات صلة كافية (أقل من العتبة المحددة)، يُعلم النظام المستخدم بعدم توفر معلومات كافية بدلاً من توليد إجابة غير مدعومة \EN{(Hallucination)}. \\
\bottomrule
\end{tabularx}
\end{table}

\subsubsection{UC-03: توليد تقرير مخصص}
\label{subsubsec:uc03}

\begin{table}[htbp]
\centering
\caption{الوصف التفصيلي لحالة الاستخدام UC-03: توليد تقرير مخصص.}
\label{tab:uc03}
\small
\renewcommand{\arraystretch}{1.45}
\begin{tabularx}{\textwidth}{>{\raggedright\arraybackslash}p{3.1cm} X}
\toprule
\textbf{الحقل} & \textbf{الوصف} \\
\midrule
المعرّف والاسم & \EN{UC-03 — Generate Report} \\
\midrule
الهدف & إتاحة طلب تقرير تحليلي مخصص عن بيانات المؤسسة مع إمكانية تحديد نوع المخرجات واكتشاف القيم الشاذة وتصدير التقرير بصيغ متعددة \EN{(PDF, CSV, Excel)}. \\
\midrule
الفاعل الرئيسي & مدير المؤسسة / موظف المؤسسة. \\
\midrule
الفاعلون المساعدون & الوكيل المشرف، ووكيل المحلل، ووكيل \EN{SQL}. \\
\midrule
المحفز & يطلب المستخدم تقريراً تحليلياً بصيغة محددة من واجهة المحادثة. \\
\midrule
الشروط المسبقة & توفر بيانات كافية ومحدثة في قاعدة البيانات لتلبية الطلب. \\
\midrule
الشروط اللاحقة & عرض التقرير المنسق أو إتاحة تنزيله بالصيغة المطلوبة، وتسجيل عملية التصدير في سجل التدقيق. \\
\midrule
المسار الأساسي & 1) يطلب المستخدم تقريراً بالعربية.\\ 2) يوجه الوكيل المشرف الطلب إلى وكيل المحلل.\\ 3) يتعاون وكيل المحلل مع وكيل \EN{SQL} لاستخراج البيانات، ثم يقوم بتحليل المؤشرات واكتشاف الاتجاهات.\\ 4) يُعرض التقرير المنسق مع خيارات التصدير.\\ 5) يختار المستخدم الصيغة \EN{(PDF/CSV)} فيتم تحميل الملف. \\
\midrule
التدفق البديل & \textbf{اكتشاف الشاذ:} في حال اكتشاف قيم شاذة \EN{(Outliers)} أثناء التحليل، يدرج وكيل المحلل تنبيهاً خاصاً ضمن التقرير للتنويه. \\
\bottomrule
\end{tabularx}
\end{table}

\subsubsection{UC-04: عرض الإشعارات والتنبيهات}
\label{subsubsec:uc04}

\begin{table}[htbp]
\centering
\caption{الوصف التفصيلي لحالة الاستخدام UC-04: عرض الإشعارات.}
\label{tab:uc04}
\small
\renewcommand{\arraystretch}{1.45}
\begin{tabularx}{\textwidth}{>{\raggedright\arraybackslash}p{3.1cm} X}
\toprule
\textbf{الحقل} & \textbf{الوصف} \\
\midrule
المعرّف والاسم & \EN{UC-04 — View Notifications} \\
\midrule
الهدف & إتاحة استعراض التنبيهات والإشعارات الداخلية الصادرة عن النظام مثل تحديثات المستندات أو جاهزية التقارير. \\
\midrule
الفاعل الرئيسي & موظف المؤسسة / مدير المؤسسة / مسؤول النظام. \\
\midrule
الفاعل المساعد & نظام الإشعارات الداخلية. \\
\midrule
المحفز & دخول المستخدم لوحة التحكم أو حدوث حدث يستدعي التنبيه. \\
\midrule
الشروط المسبقة & أن يكون المستخدم مسجلاً لدخوله على النظام. \\
\midrule
الشروط اللاحقة & تحديث حالة الإشعارات إلى (مقروءة) عند استعراضها. \\
\midrule
المسار الأساسي & 1) يستعرض المستخدم لوحة الإشعارات.\\ 2) تظهر قائمة بالتنبيهات مرتبة حسب التاريخ والأهمية.\\ 3) ينقر المستخدم على الإشعار للانتقال للتفاصيل المرتبطة. \\
\midrule
التدفق البديل & لا يوجد تدفقات بديلة رئيسية. \\
\bottomrule
\end{tabularx}
\end{table}

\subsubsection{UC-05: إدارة المستخدمين والأدوار}
\label{subsubsec:uc05}

\begin{table}[htbp]
\centering
\caption{الوصف التفصيلي لحالة الاستخدام UC-05: إدارة المستخدمين والأدوار.}
\label{tab:uc05}
\small
\renewcommand{\arraystretch}{1.45}
\begin{tabularx}{\textwidth}{>{\raggedright\arraybackslash}p{3.1cm} X}
\toprule
\textbf{الحقل} & \textbf{الوصف} \\
\midrule
المعرّف والاسم & \EN{UC-05 — Manage Users \& Roles} \\
\midrule
الهدف & تمكين مسؤول النظام من إضافة المستخدمين وتعديل بياناتهم وتحديد أدوارهم (موظف، مدير) ومستوى الأمان المرتبط بهم ضمن سياسة \EN{RBAC}. \\
\midrule
الفاعل الرئيسي & مسؤول النظام. \\
\midrule
الفاعل المساعد & قاعدة بيانات النظام. \\
\midrule
المحفز & الحاجة إلى إضافة مستخدم جديد أو تعديل صلاحيات مستخدم قائم. \\
\midrule
الشروط المسبقة & أن يكون المسؤول مسجَّل الدخول بدور \EN{(System Admin)} عبر لوحة تحكم الويب. \\
\midrule
الشروط اللاحقة & تحديث بيانات المستخدم أو دوره في قاعدة البيانات فوراً، وسريان الصلاحية الجديدة على الجلسات التالية. \\
\midrule
المسار الأساسي & 1) يفتح المسؤول لوحة إدارة المستخدمين.\\ 2) يضيف مستخدماً جديداً أو يختار مستخدماً قائماً.\\ 3) يحدد الدور (موظف/مدير) ومستوى الوصول للجداول والمستندات.\\ 4) يحفظ التغييرات فيُحدَّث النظام فوراً. \\
\midrule
التدفق البديل & في حال محاولة تعطيل حساب المسؤول الوحيد النشط، يمنع النظام العملية ويعرض تنبيهاً بضرورة وجود مسؤول واحد على الأقل. \\
\bottomrule
\end{tabularx}
\end{table}

\subsubsection{UC-06: إدارة مصادر البيانات}
\label{subsubsec:uc06}

\begin{table}[htbp]
\centering
\caption{الوصف التفصيلي لحالة الاستخدام UC-06: إدارة مصادر البيانات.}
\label{tab:uc06}
\small
\renewcommand{\arraystretch}{1.45}
\begin{tabularx}{\textwidth}{>{\raggedright\arraybackslash}p{3.1cm} X}
\toprule
\textbf{الحقل} & \textbf{الوصف} \\
\midrule
المعرّف والاسم & \EN{UC-06 — Manage Data Sources} \\
\midrule
الهدف & تمكين مسؤول النظام من ربط قواعد بيانات جديدة أو مستودعات مستندات، وتحديد هيكلها وتصنيف جداولها أمنياً. \\
\midrule
الفاعل الرئيسي & مسؤول النظام. \\
\midrule
الفاعل المساعد & طبقة \EN{Smart Views}. \\
\midrule
المحفز & ربط مؤسسة جديدة أو إضافة مستودع بيانات/مستندات حديث. \\
\midrule
الشروط المسبقة & توفر بيانات اتصال صحيحة لمصدر البيانات المستهدف. \\
\midrule
الشروط اللاحقة & تحديث طبقة \EN{Smart Views} وتأمين الاتصال بالمصدر الجديد. \\
\midrule
المسار الأساسي & 1) يدخل المسؤول لوحة «مصادر البيانات».\\ 2) يختار نوع المصدر ويُدخل بيانات الاتصال.\\ 3) يحدد الجداول ومستوى الأمان لكل منها.\\ 4) يضغط «ربط» فيختبر النظام الاتصال ويُحدّث المخطط آلياً. \\
\midrule
التدفق البديل & عند فشل الاتصال بمصدر البيانات، يظهر النظام رسالة تفصيلية بالسبب ويُتيح تعديل بيانات الاتصال. \\
\bottomrule
\end{tabularx}
\end{table}

\subsubsection{UC-07: مراجعة سجل التدقيق}
\label{subsubsec:uc07}

\begin{table}[htbp]
\centering
\caption{الوصف التفصيلي لحالة الاستخدام UC-07: مراجعة سجل التدقيق.}
\label{tab:uc07}
\small
\renewcommand{\arraystretch}{1.45}
\begin{tabularx}{\textwidth}{>{\raggedright\arraybackslash}p{3.1cm} X}
\toprule
\textbf{الحقل} & \textbf{الوصف} \\
\midrule
المعرّف والاسم & \EN{UC-07 — View Audit Log} \\
\midrule
الهدف & إتاحة اطلاع مسؤول النظام على سجل مفصل لكافة العمليات التي تمت في النظام بما فيها الاستعلامات ومسار الوكلاء لأغراض الرقابة. \\
\midrule
الفاعل الرئيسي & مسؤول النظام. \\
\midrule
الفاعل المساعد & وحدة سجلات التدقيق \EN{(Audit Logger)}. \\
\midrule
المحفز & طلب تدقيق أمني أو التحقق من سلوك مستخدم أو استكشاف خطأ. \\
\midrule
الشروط المسبقة & تسجيل الدخول بصلاحيات المسؤول \EN{(System Admin)}. \\
\midrule
الشروط اللاحقة & عرض السجلات المطلوبة أو تصديرها كملف خارجي. \\
\midrule
المسار الأساسي & 1) يفتح المسؤول لوحة «سجل التدقيق».\\ 2) يصفّي السجلات حسب المستخدم، التاريخ، أو نوع العملية.\\ 3) يختار سجلاً محدداً لعرض تفاصيله الكاملة (السؤال، الاستعلام، النتيجة، زمن الاستجابة). \\
\midrule
التدفق البديل & \textbf{تصدير السجل:} يستطيع المسؤول اختيار تصدير البيانات المفلترة كملف \EN{CSV}. \\
\bottomrule
\end{tabularx}
\end{table}

\subsubsection{UC-08: تشغيل التهيئة الذكية}
\label{subsubsec:uc08}

\begin{table}[htbp]
\centering
\caption{الوصف التفصيلي لحالة الاستخدام UC-08: تشغيل التهيئة الذكية.}
\label{tab:uc08}
\small
\renewcommand{\arraystretch}{1.45}
\begin{tabularx}{\textwidth}{>{\raggedright\arraybackslash}p{3.1cm} X}
\toprule
\textbf{الحقل} & \textbf{الوصف} \\
\midrule
المعرّف والاسم & \EN{UC-08 — Run Smart Onboarding} \\
\midrule
الهدف & إتاحة أتمتة عملية ربط النظام بمؤسسة جديدة (ربط قاعدة \EN{Odoo}، اكتشاف الجداول، وفهرسة المستندات) خلال زمن قياسي. \\
\midrule
الفاعل الرئيسي & مشغّل التهيئة \EN{(Onboarding Operator)}. \\
\midrule
الفاعلون المساعدون & نظام التحليل الآلي للجداول، وأداة الفهرسة. \\
\midrule
المحفز & الاشتراك الفعلي لمؤسسة جديدة في النظام. \\
\midrule
الشروط المسبقة & توفر بيانات اعتماد الوصول لقاعدة \EN{Odoo} والمستندات الرقمية للمؤسسة. \\
\midrule
الشروط اللاحقة & ربط قاعدة البيانات، وفهرسة المستندات في \EN{Qdrant}، وإظهار تقرير جاهزية النظام. \\
\midrule
المسار الأساسي & 1) يدخل المشغل لوحة التهيئة.\\ 2) يدخل بيانات اتصال \EN{Odoo} ويُرفع ملفات المستندات.\\ 3) يضغط «بدء التهيئة».\\ 4) يفحص النظام الاتصال، يبني \EN{Smart Views}، ويفهرس المستندات.\\ 5) يظهر تقرير جاهزية بالنظام مع إحصائيات المعالجة. \\
\midrule
التدفق البديل & \textbf{فشل الاتصال:} عند حدوث انقطاع شبكي أو خطأ بالصلاحيات، يعرض النظام تفاصيل الخطأ لتصحيح البيانات دون فقدان التقدّم المنجز. \\
\bottomrule
\end{tabularx}
\end{table}

\subsection{ملاحظات حول التحقق من حالات الاستخدام}
\label{subsec:use-case-validation}

تمثل الشروط المسبقة واللاحقة في الجداول السابقة أساساً لاشتقاق حالات الاختبار. فعلى سبيل المثال، يجب اختبار رفض الاستعلامات التي تتجاوز صلاحيات المستخدم في \EN{UC-01}، والتحقق من ظهور مصادر موثقة في \EN{UC-02}، وصحة ملفات التصدير في \EN{UC-03}. كما ينبغي اختبار الاحتفاظ بسجل تدقيق لكل عملية إدارية، وسلامة استئناف التهيئة بعد فشل الاتصال في \EN{UC-08}.

وتوفر هذه الحالات تغطية وظيفية متوازنة بين تفاعل المستخدم النهائي والإدارة والتكامل والرقابة. وبذلك تنتقل المواصفات من وصف عام للوظائف إلى سيناريوهات قابلة للتنفيذ والقياس في مرحلتي التصميم والاختبار.
